\documentclass{article}
\usepackage[utf8]{inputenc}



\begin{document}

\section{Bewertung der Gruppenmitglieder}

Bevor ich die Bewertung der Gruppenmitglieder vornehme, möchte ich kurz darauf hinweisen, dass aufgrund der zeitlichen Planung, ich nur in den Gruppentreffen und ein kleines bisschen darüber hinaus
mit Kolja, Felix H. und Florian in Kontakt stand. Daher ist die Bewertung hier etwas oberflächlich.

\subsection{Kolja}

Kolja hatte eine Phase im ersten Gruppentreffen, in der er mit seinen Problemen bezüglich Einrichtung seiner Programmierumgebung
beschäftigt war, was dazu führte, dass er zwischenzeitlich etwas abwesend wirkte. Allerdings wurde in dieser Phase etwas weniger Entscheidendes
besprochen, weshalb es nicht zu einem großen Problem wurde. Andererseits erlebte ich Kolja als einen sehr warmen und freundlichen Menschen. 
Ich errinere mich daran, dass er sich nach einem Gruppentreffen Zeit nahm, um mit mir über Schwierigkeiten im Studium und bei der Wohnungssuche 
zu sprechen. Kolja ist definitiv der Typ Mensch, auf den man immer zugehen kann. Er leitete am Ende des letzten Gruppentreffens eine 
mediatorische Diskussion ein, um eventuelle Meinungsverschiedenheiten zu klären, was bei allen gut ankam, ohne dass etwas zu besprechen war.

\subsection{Felix H.}

Felix H. war sehr engagiert dabei und schaffte mit seinen Erkenntnissen die Grundlage für meine Arbeit. Unsere Gruppentreffen verliefen trotz 
vorhandener Agenda etwas chaotisch ab, wobei Felix sich immer um Struktur in diesen und im Rest des Projektes bemühte. Leider konnte er sich nicht immer durchsetzen.
Man konnte Felix ansehen, wie er im Laufe des Projektes mehr und mehr vertraut mit der gesamten Materie wurde. Ich behaupte, 
dass er sich am Ende des Projektes am besten mit Neuralen Netzen auskennt. 

\subsection{Florian}

Florian war unsere Bank und erster Ansprechpartner bei allen technischen Problemen. Hochmotiviert stürme er sich auf die Einrichtung unseres Workspaces.
Manchmal wirkte es allerdings, als sei er in seiner eigenen Welt. Damit meine ich, dass er in Gruppentreffen manchmal vom Thema abschweifte,
um seine Gedanken zu teilen, die nicht unbedingt zum Thema passten. Dass er auch am Ende unsicher mit den Namen der Gruppenmitglieder 
war, bestätigt diesen Eindruck. Andererseits war er immer sehr freundlich und hilfsbereit. Er brachte einen nicht zu bändigen Elan 
mit, welcher auf mich ansteckend wirkte. Seine Arbeit ermöglichte es uns, uns mit der eigentlichen Materie zu beschäftigen. Der Workspace 
ist sehr gut organisiert und seine Notebooks sehr gut dokumentiert, was mir sicherlich einige Zeit und Frustation beim Einlesen zu Beginn meiner Arbeit gespart hat.

\subsection{Felix N.}

Da wir im selben Studiengang und Semester sind, hatten wir schon vor dem Projekt einige Berührungspunkte.
Felix N. ist ein klassischer Menschenfänger. Seine Art mit Menschen umzugehen und Menschen zu vernetzen ist wahrlich bewundernswert. 
Felix und ich standen während unserer Bearbeitung immer mal wieder in Kontakt, wobei wir uns über unseren Arbeitsfortschritt und das weitere Vorgehen austauschten.
Seine Motivation für das Projekt war nicht besonders hoch, was er mir gegenüber äußerte, aber für mich nicht erkennbar in seiner Arbeitsweise war.
Die Zuverlässigkeit, die er in seinen Aufgaben als Gruppensprecher an den Tag legte, sorgte auch dafür, dass ich mehr Zeit hatte um verschiedene Dinge auszuprobieren.
Obwohl er am kürzesten Zeit hatte, um am Neuronalen Netz zu arbeiten, konnte sein Ergebnis, dieses aller Gruppenmitglieder übertreffen. 

\end{document}




































\end{document}